\documentclass[../../../include/open-logic-section]{subfiles}

\begin{document}

\olfileid{sth}{ordinals}{iso}
\olsection{ક્રમ-એકરૂપતાઓ}

સુક્રમ \emph{કઈ રીતે} મજબૂત ખ્યાલ છે તે સમજાવવા માટે
સુક્રમોની સરખામણી કરવાની એક રીત રજૂ કરીને શરૂ કરીએ.

\begin{defn}
\emph{સુક્રમ} એવો જોડ $\tuple{A, <}$ છે કે $<$ એ $A$ને
સુક્રમિત કરે. સુક્રમો $\tuple{A, <}$ અને $\tuple{B,
\lessdot}$ \emph{ક્રમ-એકરૂપ} હોય જો અને ફક્ત જો એવું
!!a{bijection} $f \colon A \to B$ હોય કે:
$x < y$ તો અને ફક્ત તો $f(x) \lessdot  f(y)$.
આ કિસ્સામાં આપણે
$\ordeq{\tuple{A, <}}{\tuple{B, \lessdot}}$ લખીએ છીએ,
અને $f$ને \emph{ક્રમ-એકરૂપતા} કહીએ છીએ.
\end{defn}
\noindent
આગળ સંક્ષેપમાં આપણે “ક્રમ-એકરૂપતાઓ”ને બદલે “એકરૂપતાઓ”
કહીશું. સહજ રીતે, એકરૂપતાઓ એ સંરચના જાળવનારાં
!!{bijection}s છે. એકરૂપતાઓ વિશેની કેટલીક સરળ હકીકતો આ છે.

\begin{lem}\ollabel{isoscompose}
એકરૂપતાઓનાં સંયોજનો એકરૂપતાઓ છે; એટલે કે:
જો $f \colon A
\to B$ અને $g \colon B \to C$ એકરૂપતાઓ હોય,
તો $(g \circ f)
\colon A \to C$ પણ એકરૂપતા છે.
\end{lem}
\begin{prob}
	\olref[sth][ordinals][iso]{isoscompose} સાબિત કરો.
\end{prob}
\begin{proof}
સ્વાધ્યાય તરીકે છોડ્યું છે.
\end{proof}
\begin{cor}\ollabel{ordisoisequiv}
	$\ordeq{X}{Y}$ એ સામ્ય સંબંધ છે.
\end{cor}
\begin{prop}\ollabel{ordisounique}
જો $\tuple{A, <}$ અને $\tuple{B, \lessdot}$ એકરૂપ સુક્રમો હોય,
તો તેમની વચ્ચેની એકરૂપતા અનન્ય છે.
\end{prop}

\begin{proof}
$f$ અને $g$ એ $A \to B$ એકરૂપતાઓ છે એમ લો.
પરિણામ \olref[wo]{propwoinduction} વાપરીને અનુમાનથી
સાબિત કરીશું.
કોઈ $a\in A$ લો, અને (અનુમાન માટે)
$(\forall b < a)f(b) = g(b)$ ધારો. કોઈ $x \in B$ લો.

જો $x \lessdot f(a)$ હોય, તો $f^{-1}(x) < a$,
તેથી $f$ અને $g$ એકરૂપતાઓ છે એથી
$g(f^{-1}(x)) \lessdot
g(a)$. પણ $f^{-1}(x) < a$ હોવાથી, આપણી ધારણા મુજબ
$x =f(f^{-1}(x)) = g(f^{-1}(x))$.
આમ $x \lessdot g(a)$. એ જ રીતે, જો
$x \lessdot g(a)$ હોય, તો $x \lessdot f(a)$.

સામાન્ય રીતે,
$(\forall x \in B)(x \lessdot f(a) \liff x \lessdot
g(a))$. તેથી
\olref[sfr][rel][ord]{prop:extensionality-strictlinearorders}
મુજબ $f(a) = g(a)$. આમ
\olref[wo]{propwoinduction}થી
$(\forall
a \in A)f(a) = g(a)$.
\end{proof}\noindent
આથી સુક્રમો મજબૂત ખ્યાલ કેમ છે તેનો થોડો ખ્યાલ મળે છે.
પણ આ વાત આગળ સમજાવવા માટે થોડા વધુ સંકેતો ઉપયોગી થશે.

\begin{defn}
જ્યારે $\tuple{A, <}$ સુક્રમ હોય અને $a \in A$ હોય, ત્યારે
$A_a = \Setabs{x \in A}{x
< a}$ લો. $A_a$ને $A$નો ઉચિત \emph{આરંભિક ખંડ} કહીએ
છીએ (અને $A$ પોતે પણ $A$નો અનુચિત આરંભિક ખંડ છે એમ માનીએ છીએ).
$<_a$ એ આરંભિક ખંડ પર $<$નું મર્યાદન છે, એટલે કે
$\funrestrictionto{\mathord{<}}{A_a^2}$.
\end{defn}
\noindent
આ સંકેતો વડે કહી અને સાબિત કરી શકીએ છીએ કે કોઈ સુક્રમ
પોતાના કોઈ ઉચિત આરંભિક ખંડ સાથે એકરૂપ નથી.

\begin{lem}\ollabel{wellordnotinitial}
જો $\tuple{A, <}$ સુક્રમ હોય અને $a \in A$ હોય, તો
$\ordneq{\tuple{A, <}}{\tuple{A_a, <_a}}$
\end{lem}

\begin{proof}
વિરોધાભાસ માટે ધારો કે $f \colon A \to A_a$ એકરૂપતા
છે. $f$ એ દ્વિએક વિધેય છે અને $A_a \subsetneq A$,
તેથી \olref[wo]{wo:strictorder} વાપરીને
$A$માં એવો $<$ મુજબ લઘુતમ !!{element} $b \in A$ લો
કે $b \neq f(b)$. આપણે
$(\forall x \in A)(x<b \liff x < f(b))$ બતાવીશું;
પછી
\olref[sfr][rel][ord]{prop:extensionality-strictlinearorders}
મુજબ $b =
f(b)$ મળશે અને વિરોધાભાસ પૂર્ણ થશે.

ધારો કે $x < b$. $b$ની પસંદગી મુજબ $x = f(x)$.
વળી $f$ એકરૂપતા હોવાથી $f(x) <
f(b)$. તેથી $x < f(b)$.

ધારો કે $x < f(b)$. $f$ એકરૂપતા હોવાથી
$f^{-1}(x) < b$, અને તેથી $b$ની પસંદગી મુજબ
$f^{-1}(x) = x$. આમ $x < b$.
\end{proof}

આગળનું પરિણામ લગભગ એમ કહે છે કે એકરૂપતાનો
“આરંભિક ખંડ” પણ એકરૂપતા છે:

\begin{lem}\ollabel{wellordinitialsegment}
$\tuple{A, <}$ અને $\tuple{B, \lessdot}$ સુક્રમો હોય.
જો $f
\colon A \to B$ એકરૂપતા હોય અને $a \in A$ હોય,
તો $\funrestrictionto{f}{A_{a}} : A_a \to B_{f(a)}$
એકરૂપતા છે.
\end{lem}

\begin{proof}
$f$ એકરૂપતા હોવાથી:
	%Since $f$ is an isomorphism, $b < a$ iff $f(b) \lessdot f(a)$, so that 
\begin{align*}
	\funimage{f}{A_a} &= \funimage{f}{\Setabs{x \in A}{x < a}}\\
	&= \funimage{f}{\Setabs{f^{-1}(y) \in A}{f^{-1}(y) < a}} \\
	&= \Setabs{y \in B}{y \lessdot f(a)} \\
	&=B_{f(a)} 
\end{align*}
અને $\funrestrictionto{f}{A_a}$ ક્રમ જાળવે છે,
કારણ કે $f$ ક્રમ જાળવે છે.
\end{proof}

આગળનાં બે પરિણામો બતાવે છે કે સુક્રમોની હંમેશાં
સરખામણી થઈ શકે છે:

\begin{lem}\ollabel{lemordsegments}
$\tuple{A, <}$ અને $\tuple{B, \lessdot}$ સુક્રમો હોય.
જો
$\ordeq{\tuple{A_{a_1}, <_{a_1}}}{\tuple{B_{b_1}, \lessdot_{b_1}}}$
અને
$\ordeq{\tuple{A_{{a_2}}, <_{a_2}}}{\tuple{B_{{b_2}},
\lessdot_{b_2}}}$ હોય, તો
$ {a_1}  < {a_2} \text{ તો અને ફક્ત તો }{b_1} \lessdot
{b_2}$.
\end{lem}

\begin{proof}
આપણે \emph{ડાબેથી જમણી તરફ} સાબિત કરીશું;
બીજી દિશાની દલીલ એવી જ છે.
ધારો કે
$\ordeq{\tuple{A_{a_1}, <_{a_1}}}{\tuple{B_{b_1},
\lessdot_{b_1}}}$ અને
$\ordeq{\tuple{A_{{a_2}},
<_{a_2}}}{\tuple{B_{{b_2}}, \lessdot_{b_2}}}$ બંને છે,
અને $f \colon
A_{{a_2}} \to B_{{b_2}}$ આપણી એકરૂપતા છે.
$ {a_1} < {a_2}$ લો; ત્યારે
\olref{wellordinitialsegment}થી
$\ordeq{\tuple{A_{a_1}, <_{a_1}}}{\tuple{B_{f({a_1})},
\lessdot_{f({a_1})}}}$.
તેથી
$\ordeq{\tuple{B_{b_1}, \lessdot_{b_1}}}{\tuple{B_{f({a_1})},
\lessdot_{f({a_1})}}}$, અને આમ
\olref{wellordnotinitial}થી
$ {b_1} = f({a_1})$.
હવે $f$નો સહપ્રદેશ $B_{{b_2}}$ હોવાથી
$ {b_1} \lessdot {b_2}$.
\sourcecorrection{OLMD-115}{મૂળના આ પગલામાં $f$નો
“પ્રદેશ” $B_{{b_2}}$ કહેવાયો છે; પરંતુ ઉપર
$f\colon A_{{a_2}}\to B_{{b_2}}$ છે. અહીં
$B_{{b_2}}$ એ $f$નો સહપ્રદેશ (અને એકરૂપતા હોવાથી
વિસ્તાર) છે. આ પ્રકારે જ $b_1$ એ $b_2$થી
નાનો હોવાનું ફલિત થાય છે.}
\end{proof}

\begin{thm}\ollabel{thm:woalwayscomparable}
કોઈ પણ બે સુક્રમો આપેલા હોય, તો એક સુક્રમ બીજાના
કોઈ આરંભિક ખંડ (ઉચિત હોવો જરૂરી નથી) સાથે એકરૂપ હોય છે.
\end{thm}

\begin{proof}
$\tuple{A, <}$ અને $\tuple{B, \lessdot}$ સુક્રમો
હોય. પૃથક્કરણ વાપરીને લો
\[
	f = \Setabs{\tuple{a, b} \in A \times B}{
		\ordeq{\tuple{A_a, <_a}}{\tuple{B_b, \lessdot_b}}}.
\]
\olref{lemordsegments} મુજબ બધા
$\tuple{a_1, b_1}, \tuple{a_2, b_2} \in f$ માટે
$a_1 < a_2$ તો અને ફક્ત તો
$b_1 \lessdot b_2$. તેથી
$f \colon \dom{f} \to
\ran{f}$ એકરૂપતા છે.

જો $a_2 \in \dom{f}$ અને $a_1 < a_2$ હોય, તો
\olref{wellordinitialsegment}થી
$a_1 \in \dom{f}$; એટલે $\dom{f}$ એ $A$નો
આરંભિક ખંડ છે. એ જ રીતે, $\ran{f}$ એ $B$નો
આરંભિક ખંડ છે. વિરોધાભાસ માટે ધારો કે બંને
\emph{ઉચિત} આરંભિક ખંડ છે. ત્યારે $a$ એ
$A \setminus \dom{f}$નો $<$ મુજબ લઘુતમ
!!{element} હોય એવું લો, જેથી $\dom{f} =
A_a$; અને $b$ એ
$B \setminus
\ran{f}$નો $\lessdot$ મુજબ લઘુતમ
!!{element} હોય એવું લો, જેથી
$\ran{f} = B_b$. આમ
$f \colon A_a \to B_b$ એકરૂપતા છે અને તેથી
$\tuple{a, b} \in f$, જે વિરોધાભાસ છે.
\end{proof}

\end{document}
